% Part: computability
% Chapter: machines-computations
% Section: halting-problem

\documentclass[../../../include/open-logic-section]{subfiles}

\begin{document}

\olfileid[pt]{tur}{und}{hal}
\olsection{O Problema da Parada}

\begin{explain}
Suponha que tenhamos fixado alguma enumeração das descrições de
máquinas de Turing. Cada máquina de Turing recebe assim um
\emph{índice}: seu lugar na enumeração $M_1$, $M_2$, $M_3$, \dots{} das
descrições de máquinas de Turing.

Sabemos que devem existir funções não-Turing-computáveis: o conjunto
de descrições de máquinas de Turing---e, portanto, o conjunto de
máquinas de Turing---é !!{enumerable}, mas o conjunto de todas as
funções de $\Nat$ a $\Nat$ não é. Mas também podemos encontrar
exemplos específicos de funções não computáveis. Uma dessas funções é
a função de parada.
\end{explain}

\begin{defn}[Função de parada]
 A \emph{função de parada}~$h$ é definida como
\[
h(e,n) =
\begin{cases}
  \text{0} & \text{se a máquina~$M_e$ não para na entrada $n$} \\
  \text{1} & \text{se a máquina~$M_e$ para na entrada $n$}
\end{cases}
\]
\end{defn}

\begin{defn}[Problema da parada]
O \emph{Problema da Parada} é o problema de determinar (para quaisquer
$e$, $n$) se a máquina de Turing~$M_e$ para com uma entrada de~$n$
traços.
\end{defn}

\begin{explain}
Mostramos que~$h$ não é Turing-computável mostrando que uma função
relacionada,~$s$, não é Turing-computável. Esta prova baseia-se no
fato de que tudo o que pode ser computado por uma máquina de Turing
pode ser computado por uma máquina de Turing disciplinada
(\olref[mac][dis]{sec}), e no fato de que duas máquinas de Turing podem
ser conectadas para criar uma única máquina (\olref[mac][cmb]{sec}).
\end{explain}

\begin{defn} A função~$s$ é definida como
\[
s(e) =
\begin{cases}
  \text{0} & \text{se a máquina~$M_e$ não para na entrada $e$} \\
  \text{1} & \text{se a máquina~$M_e$ para na entrada $e$}
\end{cases}
\]
\end{defn}

\begin{lem}
A função~$s$ não é Turing-computável.
\end{lem}

\begin{proof}
Suponha, por contradição, que a função~$s$ seja Turing-computável.
Então haveria uma máquina de Turing~$S$ que computa~$s$. Podemos supor,
sem perda de generalidade, que quando $S$ para, ela o faz examinando o
primeiro quadrado (isto é, que ela é disciplinada). Esta máquina pode
ser ``conectada'' a outra máquina~$J$, que para se for iniciada na
entrada~$0$ (isto é, se ela lê $\TMblank$ no estado inicial enquanto
examina o quadrado à direita do símbolo de fim de fita), e caso
contrário vagueia para a direita, sem nunca parar. $S \concat J$, a
máquina criada conectando-se $S$ a~$J$, é uma máquina de Turing, logo é
$M_e$ para algum~$e$ (isto é, aparece em algum lugar da enumeração).
Inicie $M_e$ em uma entrada de~$e$ $\TMstroke$. Há duas
possibilidades: ou $M_e$ para ou não para.
\begin{enumerate}
\item Suponha que $M_e$ pare para uma entrada de $e$ $\TMstroke$. Então
  $s(e) = 1$. Assim, $S$, quando iniciada em~$e$, para com um único
  $\TMstroke$ como saída na fita. Então $J$ começa com um $\TMstroke$
  na fita. Nesse caso, $J$ não para. Mas $M_e$ é a máquina $S \concat
  J$, de modo que ela deveria fazer exatamente o que $S$ seguida de
  $J$ faria (isto é, neste caso, vaguear para a direita e nunca
  parar). Logo, $M_e$ não pode parar para uma entrada de $e$
  $\TMstroke$.

\item Agora suponha que $M_e$ não pare para uma entrada de $e$
  $\TMstroke$. Então $s(e) = 0$, e $S$, quando iniciada na entrada~$e$,
  para com uma fita em branco. $J$,~quando iniciada em uma fita em
  branco, para imediatamente. Novamente, $M_e$ faz o que $S$ seguida
  de~$J$ faria, de modo que $M_e$ deve parar para uma entrada de $e$
  $\TMstroke$.
\end{enumerate}
Em cada caso chegamos a uma contradição com nossa suposição. Isso
mostra que não pode haver uma máquina de Turing~$S$: $s$~não é
Turing-computável.
\end{proof}

\begin{thm}[Insolubilidade do Problema da Parada]
\ollabel{thm:halting-problem} O problema da parada é insolúvel, isto
é, a função~$h$ não é Turing-computável.
\end{thm}

\begin{proof}
Suponha que $h$ fosse Turing-computável, digamos, por uma máquina de
Turing~$H$. Poderíamos usar $H$ para construir uma máquina de Turing
que computa~$s$: primeiro, fazer uma cópia da entrada (separada por um
símbolo~$\TMblank$). Depois, voltar ao início e executar~$H$.
Claramente podemos fazer uma máquina que realiza a primeira tarefa
(veja \cref{tur:mac:dis:prob:copier}), e se $H$ existisse, poderíamos
``conectá-la'' a tal máquina copiadora para obter uma nova máquina que
determinaria se $M_e$ para na entrada~$e$, isto é, computa~$s$. Mas já
mostramos que não pode existir tal máquina. Portanto, $h$~também não é
Turing-computável.
\end{proof}

\begin{prob}
O problema da Parada com Três (3-Parada) é o problema de dar um
procedimento de decisão para determinar se uma máquina de Turing
escolhida arbitrariamente para ou não para uma entrada de três
$\TMstroke$ em uma fita que, fora isso, está em branco. Prove que o
problema 3-Parada é insolúvel.
\end{prob}

\begin{prob}
Mostre que se o problema da parada é solúvel para pares de máquina de
Turing e entrada $M_e$ e~$n$ em que $e \neq n$, então ele também é
solúvel para os casos em que $e = n$.
\end{prob}

\begin{prob}
Provamos que o problema da parada é insolúvel se a entrada é um
número~$e$, que identifica uma máquina de Turing~$M_e$ por meio de uma
enumeração de todas as máquinas de Turing. E se permitirmos a
descrição de máquinas de Turing de \olref[tur][und][enu]{sec}
diretamente como entrada? Pode haver uma máquina de Turing que decide
o problema da parada, mas que recebe como entrada descrições de
máquinas de Turing em vez de índices? Explique por que sim ou por que
não.
\end{prob}

\begin{prob} Mostre que a função \emph{parcial}~$s'$ definida como
  \[
  s'(e) =
  \begin{cases}
    \text{1} & \text{se a máquina~$M_e$ para na entrada $e$}\\
    \text{indefinida} & \text{se a máquina~$M_e$ não para na entrada $e$}
  \end{cases}
  \]
  \emph{é} Turing-computável.
\end{prob}

\end{document}
